\section{Recommender Systems}

\subsection{Collaborative Filtering}

Collaborative filtering (CF) remains a core paradigm in recommender systems, operating on the premise that user preferences can be inferred from historical user-item interactions \parencite{11359174}. Given a user set \(\mathcal{U}\), an item set \(\mathcal{I}\), and an interaction matrix \(R \in \mathbb{R}^{|\mathcal{U}| \times |\mathcal{I}|}\), the objective is to learn a scoring function
\begin{equation}
    f(u,i) \rightarrow \hat{r}_{ui},
\end{equation}
where \(\hat{r}_{ui}\) denotes the predicted relevance of item \(i \in \mathcal{I}\) to user \(u \in \mathcal{U}\). The matrix \(R\) is constructed from either explicit feedback (e.g., numerical ratings) or implicit feedback (e.g., clicks, views, purchases). In e-commerce and fashion recommendation, implicit feedback is more prevalent, representing user decisions as sparse behavioral logs and formulating recommendation as a top-\(N\) ranking task rather than rating regression.

Traditional CF methods are broadly categorized into memory-based and model-based approaches. Memory-based methods compute similarity (e.g., via cosine similarity, Jaccard coefficient, or Pearson correlation) between users or items directly from the interaction matrix. Although interpretable, their reliance on co-occurrence evidence makes them highly sensitive to data sparsity. Conversely, model-based methods learn low-dimensional latent embeddings to generalize beyond observed interactions. Matrix Factorization (MF) is the representative technique, approximating the interaction matrix as the product of user and item latent matrices:
\begin{equation}
    R \approx U V^{T},
\end{equation}
where \(U \in \mathbb{R}^{|\mathcal{U}| \times d}\) contains user latent factors, \(V \in \mathbb{R}^{|\mathcal{I}| \times d}\) contains item latent factors, and \(d\) is the embedding dimension. The relevance score is predicted via the inner product \(\hat{r}_{ui} = \mathbf{u}_{u}^{T}\mathbf{v}_{i}\).

Despite their effectiveness, traditional latent factor models suffer from several key limitations:
\begin{enumerate}
    \item \textit{Data Sparsity}: Sparse interaction logs reduce the reliability of similarity estimation and embedding quality.
    \item \textit{Cold-start Problem}: New users or items with no historical interactions cannot be effectively represented.
    \item \textit{Lack of High-order Relational Modeling}: Inner products only capture pairwise user-item correlations, ignoring multi-hop connectivity in the user-item interaction graph.
    \item \textit{Inability to Leverage Side Information}: Traditional CF cannot exploit heterogeneous features such as images or textual descriptions. This is particularly restrictive in fashion recommendation, where visual aesthetics and semantic metadata are central to user preferences.
\end{enumerate}
These limitations motivate the incorporation of content-based features, leading to alternative recommendation paradigms such as content-based filtering and hybrid systems, which are detailed in the following subsections.

\subsection{Content-based Recommendation}

Content-based (CB) recommendation systems operate on the premise that a user's future preferences can be modeled by analyzing the properties of the items they have previously interacted with. Unlike collaborative filtering, which relies on the behaviors of a wider community, CB recommendation focuses on the intrinsic characteristics of the items themselves, recommending candidate items that share high similarity with the target user's preferred catalog.

The workflow of a typical content-based recommender system consists of three main phases:
\begin{enumerate}
    \item \textit{Feature Extraction}: Items are represented as numerical feature vectors. Let \(\mathbf{x}_i \in \mathbb{R}^{d_f}\) denote the feature vector of item \(i\), which can be extracted from textual metadata (e.g., titles, descriptions using TF-IDF or transformer-based models like BERT), or visual content (e.g., clothing images processed via convolutional neural networks).
    \item \textit{Similarity Computation}: The similarity between two items \(i\) and \(j\) is measured using a distance metric over their feature vectors. A common choice is cosine similarity:
    \begin{equation}
        s(i, j) = \frac{\mathbf{x}_i^{T} \mathbf{x}_j}{\|\mathbf{x}_i\| \|\mathbf{x}_j\|},
    \end{equation}
    where \(s(i, j) \in [-1, 1]\) indicates the degree of similarity between the items.
    \item \textit{Preference Aggregation and Ranking}: For a user \(u\), a profile vector \(\mathbf{p}_u \in \mathbb{R}^{d_f}\) is constructed to represent their cumulative preferences. This is typically achieved by aggregating the feature vectors of items in their interaction history \(\mathcal{I}_u\), such as through average pooling:
    \begin{equation}
        \mathbf{p}_u = \frac{1}{|\mathcal{I}_u|} \sum_{k \in \mathcal{I}_u} \mathbf{x}_k.
    \end{equation}
    Candidate items \(j \notin \mathcal{I}_u\) are then ranked based on their similarity to the user profile, defined as \(\hat{r}_{uj} = \text{sim}(\mathbf{p}_u, \mathbf{x}_j)\), to generate the final recommendation list.
\end{enumerate}

Content-based methods offer two distinct advantages. First, they effectively mitigate the item cold-start problem: a newly introduced item can be immediately recommended as long as its descriptive features are available, without requiring historical interaction logs. Second, because recommendations are computed independently of other users' behaviors, the system is immune to the popularity bias that often degrades CF performance.

However, CB systems are constrained by two primary limitations. First, the quality of recommendations is strictly bounded by the expressiveness of the hand-crafted or pre-extracted features, making it difficult to capture complex, latent preferences that are not explicitly documented. Second, CB systems suffer from over-specialization (i.e., lack of diversity), as they tend to recommend items that are highly similar to those already consumed, thereby confining the user to a narrow information bubble and failing to offer serendipitous suggestions.


\subsection{Hybrid Recommendation Systems}

To leverage the complementary strengths of collaborative filtering and content-based approaches, hybrid recommender systems integrate both interaction histories and item side information \parencite{DBLP:journals/corr/abs-2202-02757}. By combining collaborative signals with content metadata, hybrid systems aim to construct more robust representations and produce more accurate, diverse recommendations.

Common integration strategies in hybrid systems include:
\begin{itemize}
    \item \textit{Feature-augmented Models}: Content features are directly incorporated into the collaborative filtering framework. For example, in matrix factorization, item latent vectors \(\mathbf{v}_i\) can be regularized or augmented by their corresponding content feature vectors \(\mathbf{x}_i\), allowing the model to learn representations that preserve both collaborative patterns and intrinsic attributes.
    \item \textit{Score-level Fusion}: Independent collaborative filtering and content-based recommendation models are trained separately, and their predicted scores are combined during inference. A widely used formulation is the weighted sum:
    \begin{equation}
        \hat{r}_{ui} = \alpha \hat{r}_{ui}^{\text{CF}} + (1 - \alpha) \hat{r}_{ui}^{\text{CB}},
    \end{equation}
    where \(\hat{r}_{ui}^{\text{CF}}\) and \(\hat{r}_{ui}^{\text{CB}}\) are the relevance scores predicted by the CF and CB models, respectively, and \(\alpha \in [0, 1]\) is a hyperparameter balancing the two components.
\end{itemize}

Hybrid systems are particularly valuable in modern e-commerce and fashion domains. By blending these paradigms, they effectively alleviate the data sparsity and user/item cold-start challenges that plague pure collaborative filtering. Furthermore, hybrid architectures enable the integration of multimodal data (such as product images and text descriptions), which provides rich context and significantly improves recommendation quality when interaction data is scarce \parencite{10.1007/978-981-97-9616-8_3}.

While hybrid systems successfully merge collaborative and content signals, conventional architectures often fail to fully exploit the underlying graph structure of user-item interactions or the deep visual and semantic relationships in multimodal data. This limitation motivates the adoption of graph neural networks and advanced multimodal integration strategies, which are discussed in Sections 2.2 and 2.3, respectively.
